\section{Text Elements}

    \begin{frame}
        \frametitle{Styles}
        This Slide shows different styles that Text can have. For example \textbf{bold} \textit{italic} \underline{underlined} \texttt{monospace} or \sout{canceled out}.
        
        Furthermore Text can be differently sized:
        \Huge{Huge}
        \huge{huge}
        \LARGE{LARGE}
        \large{large}
        \normalsize{normalsized}
        \small{small}
        \footnotesize{footnotesized}
        \scriptsize{scriptsize}
        \tiny{tiny}
    \end{frame}
    
    \begin{frame}
        \frametitle{Enumerations and Itemizing}
        \begin{itemize}
            \item{This is an item}
            \begin{itemize}
                \item{This is a subitem}
            \end{itemize}
        \end{itemize}
        \begin{enumerate}
            \item{First layer in enumeration}
            \begin{enumerate}
                \item{Second layer of enumeration}
            \end{enumerate}
        \end{enumerate}
    \end{frame}
    
    \begin{frame}
        \frametitle{Links to Websites}
        A link to a website can be implemented in two different ways:
        \begin{itemize}
            \item{with full link representation}
            \begin{itemize}
                \item{\href{https://www.h-brs.de/de}{https://www.h-brs.de/de}}
            \end{itemize}
            \item{with an alternative representer}
            \begin{itemize}
                \item{\href{https://www.h-brs.de/de}{Hochschule Bonn-Rhein-Sieg}}
            \end{itemize}
        \end{itemize}
    \end{frame}

    \begin{frame}
        \frametitle{Footnotes}
        Using footnotes in presentations is unusal but can be useful for more extensive explanations\footnote{For example if you are citing multiple resources in your footnotes or you do not want to overload your slide with fulltext}.
    \end{frame}
    